\documentclass[11pt]{article}
\usepackage[margin=1in]{geometry}
\usepackage{amsmath,amssymb}
\usepackage{graphicx}
\usepackage{booktabs}
\usepackage{hyperref}
\usepackage{xcolor}
\title{Replication report: \\
\emph{Phase diagram of the chiral SU(3) antiferromagnet on the kagome lattice}\\
\normalsize Xu, Capponi, Chen, Vanderstraeten, Hasik, Nevidomskyy, Mambrini, Penc, Poilblanc --- arXiv:2306.16192 (2023)}
\author{Textures-100 overnight replication (loop-current class)}
\date{\today}
\begin{document}
\maketitle

\section{Scope and classification (honest)}
This paper was filed under the \textbf{loop-current} class. That is a
\emph{partial / adjacent} fit, and we flag it explicitly. The paper is a
\emph{many-body} SU(3) \emph{spin} model --- a chiral / topological spin liquid on
the kagome lattice --- solved by exact diagonalization (Lanczos, 21/27-site tori),
matrix product states on cylinders, and PEPS/PESS tensor networks. It is not a
tight-binding loop-current \emph{metal}: there is no orbital-current order parameter
and no Peierls-flux band Chern insulator as the headline result.

The genuine loop-current connection is nonetheless real and structural. The model
\begin{equation}
H = J\!\sum_{\langle ij\rangle}\! P_{ij}
 + K_R\!\sum_{\triangle}\!\big(P_{ijk}+P_{ijk}^{-1}\big)
 + iK_I\!\sum_{\triangle}\!\big(P_{ijk}-P_{ijk}^{-1}\big)
\label{eq:H}
\end{equation}
contains a chiral three-site permutation term $iK_I$ that breaks time reversal and
reflection while preserving their product --- the defining symmetry of loop-current
order. In the single-magnon sector this term enters the $3\times3$ kagome Bloch
matrix (Eq.~A1 of the paper) \emph{exactly as an imaginary, Peierls-like hopping}
$\pm iK_I$: structurally the same time-reversal-breaking complex-hopping object that
the shared loop-current kagome kernel builds. We therefore reuse that kernel's
kagome geometry, hexagonal-BZ sampling, and batched-diagonalization machinery
(see \texttt{PROVENANCE.md}) and specialise its $3\times3$ matrix to Eq.~A1.

\paragraph{In-scope target.} The topological (TSL/CSL/double-Chern-Simons) claims
require SU(3)-symmetric tensor networks beyond an overnight single-node run and are
marked out-of-scope (not faked; see \S\ref{sec:oos}). What is fully analytical and
machine-checkable is the \textbf{single-magnon core} above the SU(3) ferromagnet,
which we replicate quantitatively.

\section{Method}
Using the sphere parametrization $J=\cos\theta\cos\phi$, $K_R=\cos\theta\sin\phi$,
$K_I=\sin\theta$ (Eq.~2), we reconstruct the single-magnon Bloch matrix Eq.~A1
(energy measured from the ferromagnet, $q=(q_x,q_y)$):
\begin{equation}
M(\mathbf q)=\!\begin{pmatrix}
-4x & 2z_-\cos\frac{q_x-\sqrt3 q_y}{2} & 2z_+\cos\frac{q_x+\sqrt3 q_y}{2}\\[2pt]
2z_+\cos\frac{q_x-\sqrt3 q_y}{2} & -4x & 2z_-\cos q_x\\[2pt]
2z_-\cos\frac{q_x+\sqrt3 q_y}{2} & 2z_+\cos q_x & -4x
\end{pmatrix},
\end{equation}
with $x\equiv J+K_R$ and $z_\pm\equiv x\pm iK_I$. We diagonalise $M(\mathbf q)$ on a
vectorised hexagonal-BZ grid and test five claims (code in
\texttt{code/}, outputs in \texttt{work/results.json}).

\section{Results}
All five machine-checkable claims are reproduced (Table~\ref{tab:res}).

\begin{table}[h]\centering
\begin{tabular}{clll}
\toprule
\# & Claim & Metric & Result\\
\midrule
1 & FM energy $e_F=2J+4K_R/3$ & max abs err (2000 pts) & $0.0$ --- PASS\\
2 & $q{=}0$ eigs $\{0,\,-6x\pm2\sqrt3 K_I\}$ & max abs err (3000 pts) & $1.4\times10^{-14}$ --- PASS\\
3 & instability $x<-|K_I|/\sqrt3$ (Eq.~3) & numeric-vs-analytic & $100.00\%$ (0/2454) --- PASS\\
4 & dispersion $\propto (x,K_I)$ only & spectrum dev (4 splits) & $0.0$ --- PASS\\
5 & flat $0$-band on boundary & lowband width on/off & $3\times10^{-15}$ / $1.5$ --- PASS\\
\bottomrule
\end{tabular}
\caption{Quantitative replication of the single-magnon core, Eq.~A1.}
\label{tab:res}
\end{table}

\paragraph{Claim 3 (ferromagnet phase boundary).} Scanning a $61\times41$ grid of
$(x,K_I)$ and computing the minimum magnon energy over a $90\times90$ BZ, the
numerically-detected stability region (min magnon energy $\ge0$) agrees with the
analytic inequality $x<-|K_I|/\sqrt3$ at every one of the 2454 points classified
away from a thin finite-grid band around the line.

\paragraph{Claim 5 + chiral hexagon modes.} On the boundary $x=-|K_I|/\sqrt3$ the
lowest magnon band collapses to a flat zero-energy band (width $\sim10^{-15}$),
reproducing the paper's ``$0$ energy band becomes flat'' statement and its
infinite-compressibility interpretation. Probing the flat-band eigenvector, the
inter-sublattice phases are exactly $\pm\pi/3$ --- confirming the paper's
chiral hexagon-mode amplitudes $e^{ij\pi/3}$.

\begin{figure}[h]\centering
\includegraphics[width=\textwidth]{../work/magnon_bands.png}
\caption{Single-magnon bands along $\Gamma$-M-K-$\Gamma$ from Eq.~A1 for a stable
ferromagnet, the boundary (flat zero band), and an unstable point.}
\end{figure}

\section{Out-of-scope items (marked, not faked)}\label{sec:oos}
The following are honest negatives, deferred to open questions 1--5: the
topological spin liquid and its single-vs-double chiral edge mode (CSL vs double
Chern--Simons); the TSL lobe extent and thermodynamic gap; the two-magnon spectra
and bound-state region (Fig.~13); and the iPEPS SU(3)-broken-phase energetics and
magnetization jump. These require ED/tensor-network machinery outside an overnight
single-node budget; the loop-current kernel's Chern/Kubo routines do not apply to a
spin model and were deliberately not forced.

\section{Verdict}
\textbf{REPRODUCED (in-scope analytical core).} 5/5 single-magnon claims match to
machine precision and the chiral hexagon-mode structure is confirmed. The
many-body topological results are the bulk of the paper and are out of scope.
\begin{itemize}
\item \textbf{Coverage: 6/10} --- full analytical single-magnon core, not the
headline many-body TSL physics.
\item \textbf{Agreement: 10/10} --- every in-scope claim reproduced to machine
precision.
\end{itemize}
\end{document}
